\documentclass[11pt,a4paper]{article}
\usepackage[utf8]{inputenc}
\usepackage[french]{babel}
\usepackage[T1]{fontenc}
\usepackage{lmodern}
\usepackage[margin=1.5cm]{geometry}

\begin{document}

% En-tête
\begin{center}
    {\Huge \textbf{Juliette Bernard}} \\[0.2cm]
    {\Large \textbf{Développeur Gameplay}} \\[0.1cm]
    juliette.bernard@email.com \ | \ 06 59 50 24 75
\end{center}

\vspace{0.3cm}

% Résumé / Profil
\section*{Profil Professionnel}
\noindent
Développeur Gameplay polyvalent avec 5 ans d'expérience dans la création d'expériences ludiques immersives. J'aime collaborer étroitement avec les Game Designers et les artistes techniques pour donner vie à des mécaniques de jeu fluides, intuitives et robustes, en m'assurant que l'expérience manette en main soit toujours parfaitement peaufinée.

\vspace{0.3cm}

% Compétences
\section*{Compétences Techniques}
\noindent
\textbf{Unreal Engine 5} $\bullet$ \textbf{Perforce} $\bullet$ \textbf{C++} $\bullet$ \textbf{Vulkan} $\bullet$ \textbf{Optimisation GPU} $\bullet$ \textbf{C\#} $\bullet$ \textbf{OpenGL} $\bullet$ \textbf{Git}

\vspace{0.3cm}

% Expériences
\section*{Expériences Professionnelles}
\noindent \textbf{Développeur Gameplay / Confirmé} --- \textit{Asobo Studio} \hfill \textbf{04/2021 à Aujourd'hui} \\
Intégration d'un moteur physique multi-threadé permettant de simuler des destructions de décors complexes à grande échelle sans perte de framerate. Création d'un système d'intelligence artificielle (Behavior Trees) pour les PNJ, rendant leurs réactions plus réalistes et adaptatives face au joueur. Développement d'un système de cinématiques temps réel, permettant de synchroniser les animations faciales avec le doublage audio. Intégration d'outils de télémétrie in-game pour récolter les données de parcours des joueurs et ajuster la difficulté dynamiquement. \\[0.3cm]
\noindent \textbf{Développeur Gameplay} --- \textit{Ubisoft} \hfill \textbf{10/2020 à 01/2021} \\
Optimisation de la gestion de la mémoire et réduction du temps de chargement des niveaux de 40\% grâce au streaming asynchrone des assets 3D. Développement de mécaniques de gameplay en C++ et optimisation du rendu graphique sur Unreal Engine pour un titre AAA, garantissant 60 FPS constants. Portage complet du moteur de jeu vers les consoles nouvelle génération (PS5, Xbox Series X), en exploitant au maximum les architectures SSD. Développement d'outils internes pour l'éditeur de niveau, augmentant de 30\% la productivité de l'équipe de Level Design au quotidien. \\[0.3cm]
\noindent \textbf{Stage — Assistant Développeur Gameplay} --- \textit{Asobo Studio} \hfill \textbf{12/2018 à 09/2020} \\
Implémentation de shaders personnalisés et de systèmes de particules avancés avec Niagara pour améliorer la qualité visuelle globale du jeu. Profilage CPU/GPU intensif à l'aide de PIX et RenderDoc pour identifier et corriger les goulots d'étranglement avant la phase de certification console. Refactoring complet du système d'inventaire modulaire, le rendant compatible avec la sauvegarde cloud multi-plateformes. \\[0.3cm]


\vspace{0.3cm}

% Formations
\section*{Formation \& Diplômes}
\noindent \textbf{Diplôme d'Ingénieur en Développement de Jeux Vidéo} \hfill \textbf{2017 - 2020} \\
\textit{Cnam ENJMIN} \\[0.3cm]
\noindent \textbf{Classes Préparatoires (CPGE) en Filière Scientifique} \hfill \textbf{2015 - 2017} \\
\textit{Lycée Scientifique} \\[0.3cm]


\end{document}